%==============================================================================
% FICHAMENTO CRÍTICO — Machine Learning for Prediction of Oral Cancer
% Recurrence
% Conforme NBR 6023 (referências) e NBR 10520 (citações)
%==============================================================================
\section{Fichamento: \citeonline{fatapour2023recurrence}}

% --- 1. IDENTIFICAÇÃO ---
\subsection{Identificação}
\begin{tabular}{ll}
\textbf{Título:} & Development of a machine learning model to predict
recurrence of oral tongue squamous cell carcinoma \\
\textbf{Autor(es):} & Fatapour, Yasaman; Abiri, Arash; Kuan, Edward C.;
Brody, James P. \\
\textbf{Periódico:} & Cancers \\
\textbf{Ano:} & 2023 \\
\textbf{Volume:} & 15 \\
\textbf{Número:} & 10 \\
\textbf{Páginas:} & 2769 \\
\textbf{DOI:} & \url{https://doi.org/10.3390/cancers15102769} \\
\textbf{Qualis:} & A1 (Oncologia) \\
\textbf{Tipo:} & Artigo original \\
\end{tabular}

% --- 2. OBJETIVO ---
\subsection{Objetivo}
Desenvolver um framework genérico baseado em machine learning para analisar
a base de dados SEER (Surveillance, Epidemiology, and End Results) e gerar
modelos preditivos confiáveis para recorrência do carcinoma espinocelular
de língua oral (OTSCC) em 5 e 10 anos, utilizando marcadores prognósticos
simples e comumente disponíveis na prática clínica.

% --- 3. METODOLOGIA ---
\subsection{Metodologia}
\textbf{Desenho do estudo:} Estudo retrospectivo de desenvolvimento e
validação de modelos preditivos baseados em ML a partir da base SEER
(2000--2018). \\
\textbf{Amostra:} 130.979 pacientes com OTSCC, sendo 36.042 (27,5\%) do
sexo feminino, idade média de 58,2 anos (DP 13,7); 13.873 não-recorrência
vs. 657 recorrência em 5 anos e 6.129 vs. 971 em 10 anos. \\
\textbf{Métricas principais:} Acurácia, precisão, recall (sensibilidade),
AUC-ROC, análise de importância de características por Shapley Additive
Explanation (SHAP). \\
\textbf{Validação:} Plataforma H2O AI com AutoML; 4 modelos testados
(Gradient Boosting Machine, Distributed Random Forest, Deep Learning,
Logistic Regression); 5-fold cross-validation; under-sampling para
balanceamento; 80\% treino / 20\% teste; 5 execuções independentes.

% --- 4. RESULTADOS PRINCIPAIS ---
\subsection{Resultados Principais}
\begin{itemize}
  \item O modelo Gradient Boosting Machine (GBM) apresentou o melhor
        desempenho, com acurácia de 81,8\% e precisão de 97,7\% para
        predição de recorrência em 5 anos.
  \item Para predição em 10 anos, o GBM alcançou acurácia de 80,0\%
        e precisão de 94,0\%.
  \item As características mais preditivas foram: número de tumores
        prévios, idade ao diagnóstico, sítio de recorrência e
        histologia tumoral, conforme análise SHAP.
  \item A aplicação de under-sampling para balanceamento das classes
        melhorou significativamente a sensibilidade do modelo sem
        comprometer a especificidade, demonstrando a importância do
        tratamento do desbalanceamento em bases de dados oncológicas
        de larga escala.
\end{itemize}

% --- 5. LIMITAÇÕES ---
\subsection{Limitações}
\begin{itemize}
  \item A base SEER apresenta limitações inerentes quanto à qualidade
        e completude dos dados, incluindo ausência de informações sobre
        tratamento sistêmico detalhado, comorbidades e status
        tabagista/etilista.
  \item A definição de recorrência no SEER depende de codificação
        secundária, que pode subestimar a taxa real de recorrência
        em comparação com dados de prontuário clínico.
  \item O estudo concentrou-se exclusivamente em carcinoma de língua
        oral (OTSCC), não abrangendo outros sítios de carcinoma
        espinocelular oral, o que limita a generalização dos achados.
  \item A precisão dos modelos para predição em 5 anos (97,7\%) é
        artificialmente elevada devido ao desbalanceamento residual
        mesmo após under-sampling; a acurácia de 81,8\% é uma
        estimativa mais realista do desempenho geral.
\end{itemize}

% --- 6. CONTRIBUIÇÃO PARA O LIVRO ---
\subsection{Contribuição para o Livro}
Este artigo fundamenta diretamente o Capítulo 12, que aborda machine learning
para predição de recorrência de câncer oral. O estudo é particularmente
relevante por quatro aspectos: (1) demonstra o uso de uma base de dados
populacional de larga escala (SEER) com 130.979 pacientes, estabelecendo
um padrão de referência para validação de modelos preditivos; (2) compara
sistematicamente quatro algoritmos de ML (GBM, DRF, Deep Learning, GLM),
fornecendo um benchmark para a seleção de modelos na seção de experimentos
do livro; (3) a análise SHAP de importância de características oferece uma
abordagem de explicabilidade (XAI) que se alinha ao framework SDD/TDD
proposto; (4) o tratamento do desbalanceamento com under-sampling é uma
lição metodológica importante para o pipeline de experimentação.

% --- 7. ANÁLISE CRÍTICA ---
\subsection{Análise Crítica}
\begin{itemize}
  \item \textbf{Validade interna:} Alta -- Framework robusto com 5-fold
        cross-validation, 5 execuções independentes e comparação
        sistemática de múltiplos algoritmos com AutoML.
  \item \textbf{Validade externa:} Média -- Dados SEER representam
        aproximadamente 28\% da população dos EUA, com diversidade
        étnica e geográfica, mas a validação em coortes internacionais
        não foi realizada.
  \item \textbf{Reprodutibilidade:} Alta -- Código disponível em
        plataforma H2O.ai; metodologia detalhada com pipeline de
        extração de dados SEER documentado passo a passo.
  \item \textbf{Relevância clínica:} Alta -- A capacidade de predizer
        recorrência com acurácia de 81,8\% pode auxiliar na
        estratificação de risco e personalização do seguimento
        oncológico pós-tratamento.
  \item \textbf{Conflito de interesses:} Não -- Os autores declararam
        ausência de conflitos de interesse.
  \item \textbf{Viés identificado:} Viés de classificação — a definição
        de recorrência baseada em codificação secundária do SEER pode
        diferir da definição clínica padrão (RECIST), introduzindo
        viés de mensuração.
\end{itemize}

% --- 8. CITAÇÕES RELEVANTES ---
\subsection{Citações Relevantes}
\begin{citacao}
``The Gradient Boosting Machine model performed the best, achieving 81.8\%
accuracy and 97.7\% precision for 5-year prediction. Moreover, 10-year
predictions demonstrated 80.0\% accuracy and 94.0\% precision.'' (p. 8).
\end{citacao}

\begin{citacao}
``The number of prior tumors, patient age, the site of cancer recurrence,
and tumor histology were the most significant predictors.'' (p. 9).
\end{citacao}

% --- 9. MAPEAMENTO SDD ---
\subsection{Mapeamento SDD}
\textbf{Problema:} Predição de recorrência locorregional do carcinoma
espinocelular de língua oral em janelas de 5 e 10 anos para estratificação
de risco e personalização do seguimento oncológico. \\
\textbf{Entrada:} Dados demográficos e clínicos extraídos do SEER (idade,
sexo, raça, estado civil, ano do diagnóstico, número de tumores prévios,
sítio tumoral, histologia, estadiamento). \\
\textbf{Saída:} Probabilidade de recorrência em 5 anos e 10 anos
(classificação binária: recorrente vs. não-recorrente). \\
\textbf{Critérios:} Acurácia $\geq$ 80\%; precisão $\geq$ 90\%;
AUC-ROC $\geq$ 0,80; identificação das top-5 características preditivas
com SHAP. \\
\textbf{Confiança:} 0,83 — Framework metodologicamente sólido com
validação cruzada e análise SHAP; limitado pela natureza retrospectiva
dos dados SEER e pela restrição a OTSCC.
\end{tabular}

% --- 10. REFERÊNCIA ABNT ---
\subsection{Referência ABNT}
\begin{verbatim}
FATAPOUR, Yasaman et al. Development of a machine learning model to predict
recurrence of oral tongue squamous cell carcinoma. Cancers, Basel, v. 15,
n. 10, p. 2769, 2023. DOI: 10.3390/cancers15102769.
\end{verbatim}
